% Generated by report/tools/tables.py from the result files of the repository. Do not edit.
\begin{table}[tbp]
\centering
\footnotesize
\setlength{\tabcolsep}{4pt}
\caption[E1 : trente essais du même programme]{Expérience E1 : trente entraînements de PPO PyTorch (Stable-Baselines3) sur HalfCheetah-v5, même code et mêmes hyperparamètres, graines aléatoires 0 à 29 ; score final : moyenne de 10 épisodes d'évaluation sur des graines jamais vues à l'entraînement. Partie haute : statistiques des 30 scores ; le test de Shapiro-Wilk rejette la normalité (trois essais décollent au-delà de 3\,000). Partie médiane : 10\,000 tirages par ligne (générateur \texttt{default\_rng(2018)}), test de Welch sur les scores finaux de deux groupes du même programme : toute différence déclarée est un faux positif. L'écart apparent divise la différence des deux moyennes par la plus petite des deux. Les groupes tirés au hasard trompent le test moins souvent que le seuil de 5~\% ; \textbf{en gras}, les groupes triés (5 meilleurs contre 5 pires sur 10 essais tirés), qui le trompent près d'une fois sur trois. Partie basse : moyennes sur 10\,000 tirages de 10 essais. Source : \file{sim/analysis/e1.json} (\file{sim/analysis/e1.py}).}
\label{tab:e1}
\begin{tabularx}{\linewidth}{@{}>{\raggedright\arraybackslash}Xrrr@{}}
\toprule
\multicolumn{4}{@{}l}{PPO PyTorch, HalfCheetah-v5, $n = 30$} \\
\midrule
\multicolumn{4}{@{}l}{\emph{Scores finaux des 30 essais (évaluation à 1 million de pas)}} \\
Moyenne ; écart type & 1\,688 ; 568 & Coef. de variation & 34~\% \\
Médiane & 1\,580 & Minimum ; maximum & 1\,138 ; 3\,430 \\
Shapiro-Wilk, $p$ & 4\,×\,10\textsuperscript{−7} & IQM & voir tableau~\ref{tab:e2} \\
\midrule
Groupes comparés & \makecell[r]{Faux positifs\\{}($p<0{,}05$)} & \makecell[r]{Différence apparent\\{}médian} & \makecell[r]{Différence dépassé\\{}1 fois sur 20} \\
\cmidrule(lr){1-4}
3 contre 3, au hasard & 1,6~\% & 17~\% & 62~\% \\
5 contre 5, au hasard & 1,2~\% & 17~\% & 47~\% \\
10 contre 10, au hasard & 1,7~\% & 11~\% & 33~\% \\
\textbf{5 meilleurs contre 5 pires, sur 10} & \textbf{31,0~\%} & 42~\% & \makecell[r]{non calculé\\{}($p$ médian 0,09)} \\
\midrule
Garder les 3 meilleurs sur 10 & \makecell[r]{moyenne des 3\\{}2\,249} & \makecell[r]{moyenne des 10\\{}1\,692} & \makecell[r]{gonflement\\{}33~\%} \\
\bottomrule
\end{tabularx}
\end{table}
